\documentclass[10pt,a4paper]{amsart}
\usepackage[margin=19mm]{geometry}
\usepackage{amsmath,amssymb,mathtools}
\usepackage[scr=rsfs,cal=euler]{mathalfa}
\usepackage{xfrac}
\usepackage[unicode,hidelinks,bookmarks=false]{hyperref}
\newcommand{\ie}{i.e.\ }
\newcommand{\cf}{cf.\ }
\newcommand{\resp}{respectively\ }
\newcommand{\Subsection}{\subsection}
\providecommand{\nobreakdash}{\nobreak\hbox{-}}
\setlength{\emergencystretch}{2em}
\multlinegap0pt
\usepackage{bm}
\usepackage{bbm}
\usepackage{dsfont}
\usepackage{xspace}
\usepackage{cancel}
\usepackage{mathtools}
\usepackage{amsthm}
\makeatletter
\@ifundefined{proposition}{\newtheorem{proposition}{Proposition}}{}
\makeatother
\DeclareMathOperator{\Tr}{Tr}
\title[General Relativity via differential forms -- explorations in]{General Relativity via differential forms -- explorations in Plebanski's Formalism for GR}
\author{Adam Shaw}
\date{Source: arXiv:2604.20772v1 (CC BY 4.0)}
\begin{document}
\maketitle

\noindent\emph{Source and licence.} General Relativity via differential forms -- explorations in Plebanski's Formalism for GR. Version arXiv:2604.20772v1 (CC BY 4.0), distributed under the Creative Commons Attribution 4.0 International licence, as recorded in the arXiv licence metadata for this exact version and shown on the abstract page https://arxiv.org/abs/2604.20772v1. Full rights notice: licenses/arxiv-2604.20772-CC-BY-4.0.txt.\par\smallskip

\noindent\textit{Editorial scope.} Counted unit: the Proposition whose source label is \texttt{prop:ConfKahlerKillingVector} in the article (source0.tex lines 4117--4119, source label \texttt{prop:ConfKahlerKillingVector}), delivered together with the article's own complete original proof of it (source0.tex lines 4120--4185, one \texttt{proof} environment of 66 lines). No mathematics is written, repaired or re-derived: statement and proof are the article's own text line for line. Required context retained uncounted: Einstein (lines 4021--4023); d-Omega-Sigma-1 (lines 4079--4081). Every definition, display and label of those lines is retained unchanged. Omitted from this package and not redistributed: 1--4020, 4024--4078, 4082--4116, 4186--7226. Nothing inside a retained excerpt is omitted or altered: every retained body is delivered exactly as the article's lines are written, so no omission receipt of any kind accompanies this package. Citations: the retained excerpts carry no citation command at all, so no bibliography entry of the article is transcribed and no reference is redistributed. Reference adaptation: none was needed and none was made; every reference inside the delivered unit resolves inside it, so no unresolved reference or citation is produced. Printed slips kept literal: no printed word, symbol, hypothesis or number of the counted statement or of its proof was altered. Layout and class: the article's own class is \texttt{report} with packages including amsmath, amsthm, amsfonts, accents, subfiles, tikz-cd, parskip, subcaption, float, xcolor, enumitem, changepage, ulem, hyperref; the wrapper is the lane's pinned \texttt{amsart} editorial wrapper at 10pt on A4, which restates the theorem-like environments the retained text needs and adds the article's own macro definitions that the retained text needs (\texttt{\textbackslash Tr}), with extra wrapper packages (bm, bbm, dsfont, xspace, cancel, mathtools). The delivered numbering is the unit's own. Assets: no retained span contains a figure environment, a graphics-inclusion command, a PDF-page inclusion, a table or a TikZ picture; no image file is copied into this package, the package declares no asset dependency and no image file is redistributed with it. Licence: version arXiv:2604.20772v1 of this article is distributed under the Creative Commons Attribution 4.0 International licence, as recorded in the arXiv licence metadata for this exact version and shown on the abstract page https://arxiv.org/abs/2604.20772v1. A full rights notice is delivered at licenses/arxiv-2604.20772-CC-BY-4.0.txt. Characters: every retained excerpt is pure ASCII, so the delivered engine, XeLaTeX with its default Latin Modern text font, reports no missing character.
    \begin{align}\label{Einstein}
        F^i = \left( \Psi^{ij} + \frac{\Lambda}{3} \delta^{ij} \right)\Sigma^j.
    \end{align}

    \begin{align}\label{d-Omega-Sigma-1}
    d\alpha (\Sigma^2+i \Sigma^3) + 3i \alpha (A^2+i A^3) \Sigma^1=0.
    \end{align}

\begin{proposition} \label{prop:ConfKahlerKillingVector}
    Let $(M,\Sigma^i)$ be an Einstein manifold with type D self-dual Weyl tensor. Let $\Sigma^i$ be a basis for the triple of self-dual 2-forms chosen so that the SD part of the Weyl curvature is diagonal, and $\Sigma^1$ is in the direction of the special (non-repeated) eigenvalue. Then $X^\mu = \Sigma^{1\mu \nu} \nabla_\nu \alpha^{-\frac{1}{3}}$ is a Killing vector field. Here $\alpha$ with $\alpha^2 = \frac{1}{6}\Tr{\Psi^2}$ is the repeated eigenvalue of the self-dual Weyl tensor.
\end{proposition}

\begin{proof}
    Inserting $X$ into $\Sigma^1$ we find
    \begin{align}
        (\iota_X \Sigma^1)_\mu = X^\nu \Sigma^1_{\nu \mu} = \Sigma^{1 \nu \rho} \nabla_\rho \alpha^{-\frac{1}{3}} \Sigma^1_{\nu \mu} = \nabla_\mu \alpha^{-\frac{1}{3}} \quad \Rightarrow \quad \iota_X \Sigma^1 = d\alpha^{-\frac{1}{3}}
    \end{align}
    where we have used the quaternion algebra of the 2-forms. The other two 1-forms are computed as
    \begin{align}
    (\iota_X \Sigma^2)_\mu = \Sigma^{1\sigma\rho} \nabla_\rho \alpha^{-\frac{1}{3}} \Sigma^2_{\sigma\mu}= \Sigma^1_\mu{}^\sigma \Sigma^2_\sigma{}^\rho \nabla_\rho \alpha^{-\frac{1}{3}},
      \end{align}
      where we have used the algebra of $\Sigma$'s to anti-commute $\Sigma^{1,2}$. We similarly have
     \begin{align}
     (\iota_X \Sigma^3)_\mu = \Sigma^1_\mu{}^\sigma \Sigma^3_\sigma{}^\rho \nabla_\rho \alpha^{-\frac{1}{3}}.
     \end{align}
     Introducing $\Sigma^+ = \Sigma^2+ i \Sigma^3$ we have
      \begin{align}\label{i-X-Sigma-plus}
      (\iota_X \Sigma^+)_\mu = \Sigma^1_\mu{}^\sigma \Sigma^+_\sigma{}^\rho \nabla_\rho \alpha^{-\frac{1}{3}}.
      \end{align}
      We now use (\ref{d-Omega-Sigma-1}), which we rewrite as 
      \begin{align}
      d\alpha^{-1/3} \Sigma^+ = i \alpha^{-1/3} A^+ \Sigma^1,
      \end{align}
      where we introduced $A^+=A^2+ iA^3$. Taking the Hodge dual we get
       \begin{align}
      \Sigma^+_\mu{}^\rho \nabla_\rho \alpha^{-1/3}  = i \alpha^{-1/3} \Sigma^1_\mu{}^\rho A^+_\rho.
      \end{align}
      Using this in (\ref{i-X-Sigma-plus}) we get
       \begin{align}
      (\iota_X \Sigma^+)_\mu =  i \alpha^{-1/3} \Sigma^1_\mu{}^\sigma  \Sigma^1_\sigma{}^\rho A^+_\rho = -i \alpha^{-1/3}A^+_\mu.
      \end{align}
      We now summarise the above results
      \begin{align}\label{i-X-Sigma-plus*}
      \iota_X \Sigma^1 = d \alpha^{-1/3}, \qquad \iota_X \Sigma^+= -i \alpha^{-1/3}A^+, \qquad  \qquad \iota_X \Sigma^-= i \alpha^{-1/3}A^-.
 \end{align}
 
  Having computed the components of $\iota_X \Sigma^i$, we can compute its exterior covariant derivative and use the characterisation of the $\Sigma$-Killing vectors to show that $X$ is a Killing vector field. It will be convenient to rewrite the formulas for the exterior covariant derivative of $\iota_X \Sigma^i$ in terms of the introduced objects $\Sigma^\pm, A^\pm$. We have
   \begin{align}
  d_A \iota_X \Sigma^1 = d \iota_X \Sigma^1 + \frac{1}{2i} (A^- \iota_X \Sigma^+ - A^+ \iota_X \Sigma^-), \\ \nonumber
  d_A \iota_X\Sigma^+ = d  \iota_X \Sigma^+ - i A^+ \iota_X \Sigma^1 + i A^1 \iota_X \Sigma^+.
    \end{align}
    Substituting (\ref{i-X-Sigma-plus*}) into the first line we see it vanishes $d_A \iota_X \Sigma^1=0$. For the first term this is manifest from $d^2=0$. In the second term there is a cancellation of the two $A^+ A^-$ terms. On the second line we have
    \begin{align}
     d_A \iota_X\Sigma^+ = -i d( \alpha^{-1/3} A^+)  - i A^+ d \alpha^{-1/3} +  A^1 \alpha^{-1/3}A^+ = - i \alpha^{-1/3} (dA^+ +i A^1 A^+).
    \end{align}
What arises here is the curvature 
 \begin{align}
 F^+ = F^2 +i  F^3 = dA^+ +i A^1 A^+.
 \end{align}
 Using the Einstein equations in the form (\ref{Einstein}) we have
 \begin{align}
 F^+ = (-\alpha + \frac{\Lambda}{3}) \Sigma^+.
\end{align}
 This means that all in all
 \begin{align}
 d_A \iota_X\Sigma^1=0, \qquad d_A \iota_X\Sigma^+ = -i \alpha^{-1/3}(-\alpha + \frac{\Lambda}{3}) \Sigma^+.
 \end{align}
    
    On the other hand, the $\Sigma$-Killing condition with the vector $\theta^i=(\theta,0,0)$ becomes
    \begin{align}
    d_A \iota_X \Sigma^1 = 0, \qquad d_A \iota_X\Sigma^+= i \theta \Sigma^+.
     \end{align}
     We see that this is indeed satisfied, with
     \begin{align}
     \theta = \alpha^{-1/3}(\alpha - \frac{\Lambda}{3}).
     \end{align}
     This finishes the proof. 
   \end{proof}

\end{document}
